\section{Infinite products of FPSs}

\subsection{Determining coefficients in products}

\begin{definition}
\label{def.fps.determines-xn-coeff}
\lean{PowerSeries.DeterminesCoeffInProd, AlgebraicCombinatorics.FPS.DeterminesCoeff}
\leantarget
\leanok
Let $\left(\mathbf{a}_{i}\right)_{i\in I}\in K\left[\left[x\right]\right]^{I}$
be a (possibly infinite) family of FPSs. Let $n\in\mathbb{N}$. Let $M$ be a
finite subset of $I$.

\textbf{(a)} We say that $M$ \emph{determines the $x^{n}$-coefficient in the
sum of} $\left(\mathbf{a}_{i}\right)_{i\in I}$ if every finite subset $J$ of
$I$ satisfying $M\subseteq J\subseteq I$ satisfies
\[
\left[x^{n}\right]\left(\sum_{i\in J}\mathbf{a}_{i}\right)
=\left[x^{n}\right]\left(\sum_{i\in M}\mathbf{a}_{i}\right).
\]

\textbf{(b)} We say that $M$ \emph{determines the $x^{n}$-coefficient in the
product of} $\left(\mathbf{a}_{i}\right)_{i\in I}$ if every finite subset $J$
of $I$ satisfying $M\subseteq J\subseteq I$ satisfies
\[
\left[x^{n}\right]\left(\prod_{i\in J}\mathbf{a}_{i}\right)
=\left[x^{n}\right]\left(\prod_{i\in M}\mathbf{a}_{i}\right).
\]
\end{definition}

\begin{definition}
\label{def.fps.xn-coeff-fin-determined}
\lean{PowerSeries.CoeffFinitelyDeterminedInProd, AlgebraicCombinatorics.FPS.CoeffFinitelyDetermined}
\leantarget
\leanok
Let $\left(\mathbf{a}_{i}\right)_{i\in I}\in K\left[\left[x\right]\right]^{I}$
be a (possibly infinite) family of FPSs. Let $n\in\mathbb{N}$.

\textbf{(a)} We say that \emph{the $x^{n}$-coefficient in the sum of}
$\left(\mathbf{a}_{i}\right)_{i\in I}$ \emph{is finitely determined} if there
is a finite subset $M$ of $I$ that determines the $x^{n}$-coefficient in
the sum of $\left(\mathbf{a}_{i}\right)_{i\in I}$.

\textbf{(b)} We say that \emph{the $x^{n}$-coefficient in the product of}
$\left(\mathbf{a}_{i}\right)_{i\in I}$ \emph{is finitely determined} if there
is a finite subset $M$ of $I$ that determines the $x^{n}$-coefficient in
the product of $\left(\mathbf{a}_{i}\right)_{i\in I}$.
\end{definition}

\begin{proposition}
\label{prop.fps.summable=fin-det}
\lean{PowerSeries.summable_iff_coeff_finitely_determined}
\leantarget
\leanok
Let $\left(\mathbf{a}_{i}\right)_{i\in I}\in K\left[\left[x\right]\right]^{I}$
be a (possibly infinite) family of FPSs. Then:

\textbf{(a)} The family $\left(\mathbf{a}_{i}\right)_{i\in I}$ is summable if
and only if each coefficient in its sum is finitely determined (i.e., for each
$n\in\mathbb{N}$, the $x^{n}$-coefficient in the sum of
$\left(\mathbf{a}_{i}\right)_{i\in I}$ is finitely determined).

\textbf{(b)} If the family $\left(\mathbf{a}_{i}\right)_{i\in I}$ is summable,
then its sum $\sum_{i\in I}\mathbf{a}_{i}$ is the FPS whose $x^{n}$-coefficient
(for any $n\in\mathbb{N}$) can be computed as follows: If $n\in\mathbb{N}$, and
if $M$ is a finite subset of $I$ that determines the $x^{n}$-coefficient in the
sum of $\left(\mathbf{a}_{i}\right)_{i\in I}$, then
\[
\left[x^{n}\right]\left(\sum_{i\in I}\mathbf{a}_{i}\right)
=\left[x^{n}\right]\left(\sum_{i\in M}\mathbf{a}_{i}\right).
\]
\end{proposition}

\begin{proof}
Easy and LTTR.
\end{proof}

\subsection{Multipliable families and infinite products}

\begin{definition}
\label{def.fps.multipliable}
\lean{PowerSeries.Multipliable, AlgebraicCombinatorics.FPS.Multipliable}
\leantarget
\leanok
Let $\left(\mathbf{a}_{i}\right)_{i\in I}$ be a (possibly infinite) family of
FPSs. Then:

\textbf{(a)} The family $\left(\mathbf{a}_{i}\right)_{i\in I}$ is said to be
\emph{multipliable} if and only if each coefficient in its product is finitely
determined.

\textbf{(b)} If the family $\left(\mathbf{a}_{i}\right)_{i\in I}$ is
multipliable, then its \emph{product} $\prod_{i\in I}\mathbf{a}_{i}$ is defined
to be the FPS whose $x^{n}$-coefficient (for any $n\in\mathbb{N}$) can be
computed as follows: If $n\in\mathbb{N}$, and if $M$ is a finite subset of $I$
that determines the $x^{n}$-coefficient in the product of $\left(
\mathbf{a}_{i}\right)_{i\in I}$, then
\[
\left[x^{n}\right]\left(\prod_{i\in I}\mathbf{a}_{i}\right)
=\left[x^{n}\right]\left(\prod_{i\in M}\mathbf{a}_{i}\right).
\]
\end{definition}

\begin{proposition}
\label{prop.fps.multipliable.prod-wd}
\lean{PowerSeries.multipliable_coeff_eq_of_determines}
\leantarget
\leanok
This definition of $\prod_{i\in I}\mathbf{a}_{i}$ is well-defined -- i.e., the
coefficient $\left[x^{n}\right]\left(\prod_{i\in M}\mathbf{a}_{i}\right)$ does
not depend on $M$ (as long as $M$ is a finite subset of $I$ that determines the
$x^{n}$-coefficient in the product of
$\left(\mathbf{a}_{i}\right)_{i\in I}$).
\end{proposition}

\begin{proof}
Let $n\in\mathbb{N}$. We need to check that the coefficient
$\left[x^{n}\right]\left(\prod_{i\in M}\mathbf{a}_{i}\right)$ does not depend
on $M$. In other words, we need to check that if $M_{1}$ and $M_{2}$ are two
finite subsets of $I$ that each determine the $x^{n}$-coefficient in the product
of $\left(\mathbf{a}_{i}\right)_{i\in I}$, then
\[
\left[x^{n}\right]\left(\prod_{i\in M_{1}}\mathbf{a}_{i}\right)
=\left[x^{n}\right]\left(\prod_{i\in M_{2}}\mathbf{a}_{i}\right).
\]
So let us prove this. Since $M_{1}$ determines the $x^{n}$-coefficient, every
finite subset $J$ of $I$ with $M_{1}\subseteq J$ satisfies
$\left[x^{n}\right]\left(\prod_{i\in J}\mathbf{a}_{i}\right)
=\left[x^{n}\right]\left(\prod_{i\in M_{1}}\mathbf{a}_{i}\right)$.
Applying this to $J=M_{1}\cup M_{2}$, we obtain
\[
\left[x^{n}\right]\left(\prod_{i\in M_{1}\cup M_{2}}\mathbf{a}_{i}\right)
=\left[x^{n}\right]\left(\prod_{i\in M_{1}}\mathbf{a}_{i}\right).
\]
The same argument (with the roles of $M_{1}$ and $M_{2}$ swapped) yields
\[
\left[x^{n}\right]\left(\prod_{i\in M_{2}\cup M_{1}}\mathbf{a}_{i}\right)
=\left[x^{n}\right]\left(\prod_{i\in M_{2}}\mathbf{a}_{i}\right).
\]
Since $M_{1}\cup M_{2}=M_{2}\cup M_{1}$, the left hand sides are equal, whence
$\left[x^{n}\right]\left(\prod_{i\in M_{1}}\mathbf{a}_{i}\right)
=\left[x^{n}\right]\left(\prod_{i\in M_{2}}\mathbf{a}_{i}\right)$.
\end{proof}

\begin{proposition}
\label{prop.fps.multipliable.prod-wd2}
\lean{PowerSeries.tprod_eq_finprod}
\leantarget
\leanok
Let $\left(\mathbf{a}_{i}\right)_{i\in I}$ be a finite family of FPSs. Then, the
product $\prod_{i\in I}\mathbf{a}_{i}$ defined according to Definition
\ref{def.fps.multipliable} \textbf{(b)} equals the finite product
$\prod_{i\in I}\mathbf{a}_{i}$ defined in the usual way (i.e., defined as in any
commutative ring).
\end{proposition}

\begin{proof}
Argue that $I$ itself is a subset of $I$ that determines all coefficients in
the product of $\left(\mathbf{a}_{i}\right)_{i\in I}$. See Section
\ref{sec.details.gf.prod} for a detailed proof.
\end{proof}

\subsection{Irrelevance lemmas}

\begin{lemma}
\label{lem.fps.prod.irlv.cong-mul}
\lean{AlgebraicCombinatorics.FPS.xnEquiv_mul_of_coeff_eq}
\leanhelper
\leanok
If $a\overset{x^n}{\equiv} b$ and $c\overset{x^n}{\equiv} d$, then
$ac\overset{x^n}{\equiv} bd$.
In other words, $x^n$-equivalence is preserved under multiplication.
\end{lemma}

\begin{proof}
\leanok
This follows from the Cauchy product formula: the $x^m$-coefficient of a
product $ac$ depends only on the coefficients of $a$ and $c$ up to degree $m$.
\end{proof}

\begin{lemma}
\label{lem.fps.prod.irlv.cong-div}
\lean{AlgebraicCombinatorics.FPS.xnEquiv_div}
\leanhelper
\leanok
If $a\overset{x^n}{\equiv} b$ and $c\overset{x^n}{\equiv} d$ with
$c$ and $d$ invertible (i.e., $[x^0]c$ and $[x^0]d$ are units), then
$a/c \overset{x^n}{\equiv} b/d$.
\end{lemma}

\begin{proof}
\leanok
First show that $c^{-1}\overset{x^n}{\equiv} d^{-1}$ using
$x^n$-equivalence of inverses, then apply Lemma
\ref{lem.fps.prod.irlv.cong-mul} to $a\cdot c^{-1}$ and $b\cdot d^{-1}$.
\end{proof}

\begin{lemma}
\label{lem.fps.prod.irlv.1}
\lean{PowerSeries.coeff_mul_one_add_eq_of_coeff_zero}
\leantarget
\leanok
Let $a,f\in K\left[\left[x\right]\right]$ be two FPSs. Let $n\in\mathbb{N}$.
Assume that
\[
\left[x^{m}\right]f=0 \qquad\text{for each }m\in\left\{0,1,\ldots,n\right\}.
\]
Then,
\[
\left[x^{m}\right]\left(a\left(1+f\right)\right)=\left[x^{m}\right]a
\qquad\text{for each }m\in\left\{0,1,\ldots,n\right\}.
\]
\end{lemma}

\begin{proof}
The FPS $af$ is a multiple of $f$ (since $af=fa$). Hence, Lemma
\ref{lem.fps.prod.irlv.mul} yields that
$\left[x^{m}\right]\left(af\right)=0$ for each
$m\in\left\{0,1,\ldots,n\right\}$.
Now, for each $m\in\left\{0,1,\ldots,n\right\}$, we have
\[
\left[x^{m}\right]\left(a\left(1+f\right)\right)
=\left[x^{m}\right]\left(a+af\right)
=\left[x^{m}\right]a+\underbrace{\left[x^{m}\right]\left(af\right)}_{=0}
=\left[x^{m}\right]a.
\]
\end{proof}

\begin{lemma}
\label{lem.fps.prod.irlv.fin}
\lean{PowerSeries.coeff_mul_prod_one_add_eq}
\leantarget
\leanok
Let $a\in K\left[\left[x\right]\right]$ be an FPS. Let $\left(f_{i}\right)_{i\in J}
\in K\left[\left[x\right]\right]^{J}$ be a finite family of FPSs. Let
$n\in\mathbb{N}$. Assume that each $i\in J$ satisfies
\[
\left[x^{m}\right]\left(f_{i}\right)=0 \qquad\text{for each }
m\in\left\{0,1,\ldots,n\right\}.
\]
Then,
\[
\left[x^{m}\right]\left(a\prod_{i\in J}\left(1+f_{i}\right)\right)
=\left[x^{m}\right]a \qquad\text{for each }m\in\left\{0,1,\ldots,n\right\}.
\]
\end{lemma}

\begin{proof}
This is just Lemma \ref{lem.fps.prod.irlv.1}, applied $\left|J\right|$ many
times. See Section \ref{sec.details.gf.prod} for a
detailed proof.
\end{proof}

\subsection{Multipliability criterion}

\begin{theorem}
\label{thm.fps.1+f-mulable}
\lean{PowerSeries.multipliable_one_add_of_summable}
\leantarget
\leanok
Let $\left(f_{i}\right)_{i\in I}\in K\left[\left[x\right]\right]^{I}$ be a
(possibly infinite) summable family of FPSs. Then, the family
$\left(1+f_{i}\right)_{i\in I}$ is multipliable.
\end{theorem}

\begin{proof}
\leanok
This is an easy consequence of Lemma \ref{lem.fps.prod.irlv.fin}. See Section
\ref{sec.details.gf.prod} for a detailed proof.
\end{proof}

\begin{proposition}
\label{prop.fps.1-mulable}
\lean{PowerSeries.multipliable_of_finite_ne_one}
\leantarget
\leanok
If all but finitely many entries of a family
$\left(\mathbf{a}_{i}\right)_{i\in I}\in K\left[\left[x\right]\right]^{I}$
equal $1$ (that is, if all but finitely many $i\in I$ satisfy
$\mathbf{a}_{i}=1$), then this family is multipliable.
\end{proposition}

\begin{proof}
\leanok
LTTR. (See Section \ref{sec.details.gf.prod} for a detailed proof.)
\end{proof}

\subsection{Further multipliability results}

\begin{lemma}
\label{lem.fps.multipliable-const-one}
\lean{PowerSeries.multipliable_const_one}
\leanhelper
\leanok
If every entry of a family
$\left(\mathbf{a}_{i}\right)_{i\in I}$ equals $1$, then the family is multipliable.
\end{lemma}

\begin{proof}
\leanok
The empty set determines every coefficient (the product over $\emptyset$ is $1$).
\end{proof}

\begin{lemma}
\label{lem.fps.multipliable-empty}
\lean{PowerSeries.multipliable_empty}
\leanhelper
\leanok
Any family indexed by the empty type is multipliable.
\end{lemma}

\begin{proof}
\leanok
The empty set determines all coefficients.
\end{proof}

\begin{lemma}
\label{lem.fps.multipliable-singleton}
\lean{PowerSeries.multipliable_singleton}
\leanhelper
\leanok
Any family indexed by a unique (singleton) type is multipliable.
\end{lemma}

\begin{proof}
\leanok
The single-element set determines all coefficients.
\end{proof}

\begin{lemma}
\label{lem.fps.multipliable-of-zero-mem}
\lean{PowerSeries.multipliable_of_zero_mem}
\leanhelper
\leanok
If some $\mathbf{a}_j = 0$ in the family, then the family is multipliable
(the product is $0$).
\end{lemma}

\begin{proof}
\leanok
Any finite set containing $j$ determines all coefficients (the product is $0$).
\end{proof}

\begin{lemma}
\label{lem.fps.multipliable-of-finite}
\lean{AlgebraicCombinatorics.FPS.multipliable_of_finite}
\leanhelper
\leanok
If $I$ is a finite set, then any family $\left(\mathbf{a}_{i}\right)_{i\in I}$
is multipliable.
\end{lemma}

\begin{proof}
\leanok
$I$ itself determines all coefficients.
\end{proof}

\begin{lemma}
\label{lem.fps.infprod-of-finite}
\lean{AlgebraicCombinatorics.FPS.infprod_of_finite}
\leanhelper
\leanok
If $I$ is a finite set and $\left(\mathbf{a}_{i}\right)_{i\in I}$
is a family of FPSs, then the infinite product $\prod_{i\in I}\mathbf{a}_{i}$
equals the ordinary finite product.
\end{lemma}

\begin{proof}
\leanok
This is a special case of Proposition \ref{prop.fps.multipliable.prod-wd2}.
\end{proof}

\begin{lemma}
\label{lem.fps.tprod-const-one}
\lean{PowerSeries.tprod_const_one}
\leanhelper
\leanok
If every entry of a family equals $1$, then the infinite product equals $1$.
\end{lemma}

\begin{proof}
\leanok
The empty set is an $x^n$-approximator for every $n$,
and $\prod_{i\in\emptyset}\mathbf{a}_i = 1$.
\end{proof}

\subsection{Binary product example}

\begin{lemma}
\label{lem.fps.binary-prod-finite}
\lean{PowerSeries.prod_one_add_pow_two_eq}
\leanhelper
\leanok
For each $m\in\mathbb{N}$,
\[
\prod_{i=0}^{m}\left(1+x^{2^{i}}\right)
=\left(1-x^{2^{m+1}}\right)\cdot\frac{1}{1-x}.
\]
\end{lemma}

\begin{proof}
By induction on $m$, using the identity $\left(1-a\right)\left(1+a\right)=1-a^{2}$.
\end{proof}

\begin{lemma}
\label{lem.fps.binary-prod-mulable}
\lean{PowerSeries.multipliable_one_add_pow_two}
\leanhelper
\leanok
The family $\left(1+x^{2^{i}}\right)_{i\in\mathbb{N}}$ is multipliable.
\end{lemma}

\begin{proof}
\leanok
Apply Theorem \ref{thm.fps.1+f-mulable}: the family
$\left(x^{2^{i}}\right)_{i\in\mathbb{N}}$ is summable (for each $n$, at most
one $i$ satisfies $2^i = n$ since $i\mapsto 2^i$ is injective).
\end{proof}

\begin{proposition}
\label{prop.fps.binary-prod-infinite}
\lean{PowerSeries.tprod_one_add_pow_two_eq}
\leanhelper
\leanok
We have
\[
\prod_{i=0}^{\infty}\left(1+x^{2^{i}}\right) = \frac{1}{1-x}.
\]
This relates to the unique binary representation of natural numbers.
\end{proposition}

\begin{proof}
\leanok
For each $n$, the set $\{0,1,\ldots,n\}$ determines the $x^n$-coefficient:
any factor $\left(1+x^{2^i}\right)$ with $i > n$ has $2^i > n$, so
it does not affect coefficients $\le n$. The finite product for
$i\in\{0,\ldots,n\}$ equals $\left(1-x^{2^{n+1}}\right)/(1-x)$
by Lemma \ref{lem.fps.binary-prod-finite}, and its $x^n$-coefficient
equals that of $1/(1-x)$ since $2^{n+1} > n$.
\end{proof}

\subsection{$x^n$-approximators}

\begin{definition}
\label{def.fps.infprod-approx}
\lean{PowerSeries.IsXnApproximator, AlgebraicCombinatorics.FPS.IsXnApproximator}
\leantarget
\leanok
Let $\left(\mathbf{a}_{i}\right)_{i\in I}\in K\left[\left[x\right]\right]^{I}$
be a family of FPSs. Let $n\in\mathbb{N}$. An $x^{n}$\emph{-approximator} for
$\left(\mathbf{a}_{i}\right)_{i\in I}$ means a finite subset $M$ of $I$ that
determines the first $n+1$ coefficients in the product of
$\left(\mathbf{a}_{i}\right)_{i\in I}$. (In other words, $M$ has to determine
the $x^{m}$-coefficient in the product of $\left(\mathbf{a}_{i}\right)_{i\in I}$
for each $m\in\left\{0,1,\ldots,n\right\}$.)
\end{definition}

\begin{lemma}
\label{lem.fps.mulable.approx}
\lean{PowerSeries.exists_xn_approximator, AlgebraicCombinatorics.FPS.multipliable_exists_approximator}
\leantarget
\leanok
Let $\left(\mathbf{a}_{i}\right)_{i\in I}\in K\left[\left[x\right]\right]^{I}$
be a multipliable family of FPSs. Let $n\in\mathbb{N}$. Then, there exists an
$x^{n}$-approximator for $\left(\mathbf{a}_{i}\right)_{i\in I}$.
\end{lemma}

\begin{proof}
\leanok
This is an easy consequence of the fact that a union of finitely many finite
sets is finite. A detailed proof can be found in Section
\ref{sec.details.gf.prod}.
\end{proof}

\begin{proposition}
\label{prop.fps.infprod-approx-xneq}
\lean{PowerSeries.xn_approximator_superset_xnEquiv, PowerSeries.tprod_xnEquiv_approximator}
\leantarget
\leanok
Let $\left(\mathbf{a}_{i}\right)_{i\in I}\in K\left[\left[x\right]\right]^{I}$
be a family of FPSs. Let $n\in\mathbb{N}$. Let $M$ be an $x^{n}$-approximator
for $\left(\mathbf{a}_{i}\right)_{i\in I}$. Then:

\textbf{(a)} Every finite subset $J$ of $I$ satisfying $M\subseteq J\subseteq I$
satisfies
\[
\prod_{i\in J}\mathbf{a}_{i}\overset{x^{n}}{\equiv}\prod_{i\in M}\mathbf{a}_{i}.
\]

\textbf{(b)} If the family $\left(\mathbf{a}_{i}\right)_{i\in I}$ is
multipliable, then
\[
\prod_{i\in I}\mathbf{a}_{i}\overset{x^{n}}{\equiv}\prod_{i\in M}\mathbf{a}_{i}.
\]
\end{proposition}

\begin{proof}
This follows easily from Definition \ref{def.fps.infprod-approx} and
Definition \ref{def.fps.multipliable} \textbf{(b)}. See Section
\ref{sec.details.gf.prod} for a detailed proof.
\end{proof}

\subsection{Properties of infinite products}

\begin{proposition}
\label{prop.fps.union-mulable}
\lean{PowerSeries.multipliable_of_union, PowerSeries.tprod_eq_tprod_mul_tprod}
\leantarget
\leanok
Let $\left(\mathbf{a}_{i}\right)_{i\in I}\in K\left[\left[x\right]\right]^{I}$
be a family of FPSs. Let $J$ be a subset of $I$. Assume that the subfamilies
$\left(\mathbf{a}_{i}\right)_{i\in J}$ and
$\left(\mathbf{a}_{i}\right)_{i\in I\setminus J}$ are multipliable. Then:

\textbf{(a)} The entire family $\left(\mathbf{a}_{i}\right)_{i\in I}$ is
multipliable.

\textbf{(b)} We have
\[
\prod_{i\in I}\mathbf{a}_{i}=\left(\prod_{i\in J}\mathbf{a}_{i}\right)
\cdot\left(\prod_{i\in I\setminus J}\mathbf{a}_{i}\right).
\]
\end{proposition}

\begin{proof}
Fix $n\in\mathbb{N}$. Lemma \ref{lem.fps.mulable.approx} (applied to $J$
instead of $I$) shows that there exists an $x^{n}$-approximator $U$ for
$\left(\mathbf{a}_{i}\right)_{i\in J}$. Lemma \ref{lem.fps.mulable.approx}
shows that there exists an $x^{n}$-approximator $V$ for
$\left(\mathbf{a}_{i}\right)_{i\in I\setminus J}$. Note that $U\cup V$ is
finite. Now, $U\cup V$ is an $x^{n}$-approximator for
$\left(\mathbf{a}_{i}\right)_{i\in I}$. Hence, the $x^{n}$-coefficient in the
product of $\left(\mathbf{a}_{i}\right)_{i\in I}$ is finitely determined. This
proves part \textbf{(a)}. Part \textbf{(b)} follows using Proposition
\ref{prop.fps.infprod-approx-xneq} \textbf{(b)}.

See Section \ref{sec.details.gf.prod} for the details.
\end{proof}

\begin{proposition}
\label{prop.fps.prod-mulable}
\lean{PowerSeries.multipliable_mul, PowerSeries.tprod_mul_eq_mul_tprod}
\leantarget
\leanok
Let $\left(\mathbf{a}_{i}\right)_{i\in I}\in K\left[\left[x\right]\right]^{I}$
and $\left(\mathbf{b}_{i}\right)_{i\in I}\in K\left[\left[x\right]\right]^{I}$
be two multipliable families of FPSs. Then:

\textbf{(a)} The family $\left(\mathbf{a}_{i}\mathbf{b}_{i}\right)_{i\in I}$ is
multipliable.

\textbf{(b)} We have
\[
\prod_{i\in I}\left(\mathbf{a}_{i}\mathbf{b}_{i}\right)=\left(\prod_{i\in I}
\mathbf{a}_{i}\right)\cdot\left(\prod_{i\in I}\mathbf{b}_{i}\right).
\]
\end{proposition}

\begin{proof}
Fix $n\in\mathbb{N}$. Lemma \ref{lem.fps.mulable.approx} shows that there
exists an $x^{n}$-approximator $U$ for $\left(\mathbf{a}_{i}\right)_{i\in I}$.
Lemma \ref{lem.fps.mulable.approx} (applied to $\mathbf{b}_{i}$ instead of
$\mathbf{a}_{i}$) shows that there exists an $x^{n}$-approximator $V$ for
$\left(\mathbf{b}_{i}\right)_{i\in I}$. Note that $U\cup V$ is finite. Now,
$U\cup V$ is an $x^{n}$-approximator for
$\left(\mathbf{a}_{i}\mathbf{b}_{i}\right)_{i\in I}$. From here, proceed as
in the proof of Proposition \ref{prop.fps.union-mulable}.

See Section \ref{sec.details.gf.prod} for the details.
\end{proof}

\begin{proposition}
\label{prop.fps.div-mulable}
\lean{PowerSeries.multipliable_div, PowerSeries.tprod_div_eq_div_tprod}
\leantarget
\leanok
Let $\left(\mathbf{a}_{i}\right)_{i\in I}\in K\left[\left[x\right]\right]^{I}$
and $\left(\mathbf{b}_{i}\right)_{i\in I}\in K\left[\left[x\right]\right]^{I}$
be two multipliable families of FPSs. Assume that the FPS $\mathbf{b}_{i}$ is
invertible for each $i\in I$. Then:

\textbf{(a)} The family $\left(\dfrac{\mathbf{a}_{i}}{\mathbf{b}_{i}}\right)_{i\in I}$
is multipliable.

\textbf{(b)} We have
\[
\prod_{i\in I}\dfrac{\mathbf{a}_{i}}{\mathbf{b}_{i}}
=\dfrac{\prod_{i\in I}\mathbf{a}_{i}}{\prod_{i\in I}\mathbf{b}_{i}}.
\]
\end{proposition}

\begin{proof}
This is similar to the proof of Proposition \ref{prop.fps.prod-mulable}, but
using Lemma \ref{lem.fps.prod.irlv.cong-div} ($x^n$-equivalence of quotients)
instead of Lemma \ref{lem.fps.prod.irlv.cong-mul}. See Section
\ref{sec.details.gf.prod} for the details.
\end{proof}

\subsection{Subfamilies of multipliable families}

\begin{lemma}
\label{lem.fps.prods-mulable-subfams-appr}
\lean{AlgebraicCombinatorics.FPS.xnApproximator_inter}
\leanhelper
\leanok
Let $\left(\mathbf{a}_{i}\right)_{i\in I}$ be a family of invertible FPSs.
Let $U$ be an $x^n$-approximator for $\left(\mathbf{a}_{i}\right)_{i\in I}$.
Let $J\subseteq I$. Then $U\cap J$ (viewed as a subset of $J$) is an
$x^n$-approximator for the subfamily $\left(\mathbf{a}_{i}\right)_{i\in J}$.
\end{lemma}

\begin{proof}
\leanok
Let $T\supseteq U\cap J$ be a finite subset of $J$. Then
$T\cup (U\setminus J)$ and $(U\cap J)\cup (U\setminus J) = U$ are both
finite supersets of $U$, so their products are $x^n$-equivalent
to $\prod_{i\in U}\mathbf{a}_i$ by the approximator property.
Writing each product as a product over the $J$-part times a product over
$U\setminus J$, and canceling the invertible common factor
$\prod_{i\in U\setminus J}\mathbf{a}_i$, we conclude
$\prod_{i\in T}\mathbf{a}_i \overset{x^n}{\equiv} \prod_{i\in U\cap J}\mathbf{a}_i$.
\end{proof}

\begin{proposition}
\label{prop.fps.prods-mulable-subfams}
\lean{PowerSeries.multipliable_subfamily}
\leantarget
\leanok
Let $\left(\mathbf{a}_{i}\right)_{i\in I}\in K\left[\left[x\right]\right]^{I}$
be a multipliable family of invertible FPSs. Then, any subfamily of
$\left(\mathbf{a}_{i}\right)_{i\in I}$ is multipliable.
\end{proposition}

\begin{proof}
\leanok
We must show that $(\mathbf{a}_i)_{i\in J}$ is multipliable whenever
$J \subseteq I$. The idea is to show that if $U$ is an $x^{n}$-approximator
for $\left(\mathbf{a}_{i}\right)_{i\in I}$, then $U\cap J$ determines the
$x^{n}$-coefficient in the product of $\left(\mathbf{a}_{i}\right)_{i\in J}$
(and, in fact, is an $x^{n}$-approximator for
$\left(\mathbf{a}_{i}\right)_{i\in J}$). This uses Lemma
\ref{lem.fps.prods-mulable-subfams-appr} and relies on the invertibility of
$\prod_{i\in U\setminus J}\mathbf{a}_{i}$; this is why the FPS
$\mathbf{a}_{i}$ are required to be invertible in the proposition.

See Section \ref{sec.details.gf.prod} for the details.
\end{proof}

\subsection{Reindexing}

\begin{proposition}
\label{prop.fps.prods-mulable-rules.reindex}
\lean{PowerSeries.multipliable_reindex, PowerSeries.tprod_reindex}
\leantarget
\leanok
Let $S$ and $T$ be two sets. Let $f:S\rightarrow T$ be a bijection. Let
$\left(\mathbf{a}_{t}\right)_{t\in T}\in K\left[\left[x\right]\right]^{T}$
be a multipliable family of FPSs. Then,
\[
\prod_{t\in T}\mathbf{a}_{t}=\prod_{s\in S}\mathbf{a}_{f\left(s\right)}
\]
(and, in particular, the product on the right hand side is well-defined, i.e.,
the family $\left(\mathbf{a}_{f\left(s\right)}\right)_{s\in S}$ is multipliable).
\end{proposition}

\begin{proof}
This is a trivial consequence of the definitions of multipliability and
infinite products.
\end{proof}

\subsection{Breaking products into subproducts}

\begin{lemma}
\label{lem.fps.prods-mulable-rules.SW1.lem1}
\lean{AlgebraicCombinatorics.FPS.xnEquiv_finprod}
\leantarget
\leanok
Let $n\in\mathbb{N}$ and let $V$ be a finite set. Let $c,d : V \to K[[x]]$ be
two families of FPSs such that $c_v \overset{x^n}{\equiv} d_v$ for each $v\in V$.
Then
\[
\prod_{v\in V} c_v \overset{x^n}{\equiv} \prod_{v\in V} d_v.
\]
\end{lemma}

\begin{proof}
By induction on $|V|$. The base case $V = \emptyset$ is trivial. In the
inductive step, use $x^n$-equivalence of products:
$\prod_{v\in V \cup \{w\}} c_v = c_w \cdot \prod_{v\in V} c_v$
and similarly for $d$, then apply the fact that $x^n$-equivalence is preserved
under multiplication.
\end{proof}

\begin{proposition}
\label{prop.fps.prods-mulable-rules.SW1}
\lean{PowerSeries.multipliable_prod_fibers_inv, PowerSeries.tprod_eq_tprod_fibers_inv}
\leantarget
\leanok
Let $\left(\mathbf{a}_{s}\right)_{s\in S}\in K\left[\left[x\right]\right]^{S}$
be a multipliable family of FPSs. Let $W$ be a set. Let $f:S\rightarrow W$ be
a map. Assume that for each $w\in W$, the family $\left(\mathbf{a}_{s}\right)_{s\in
S,\; f(s)=w}$ is multipliable. Then,
\[
\prod_{s\in S}\mathbf{a}_{s}=\prod_{w\in W}\;\;\prod_{\substack{s\in S;\\
f(s)=w}}\mathbf{a}_{s}.
\]
(In particular, the right hand side is well-defined -- i.e., the family
$\left(\prod_{\substack{s\in S;\\ f(s)=w}}\mathbf{a}_{s}\right)_{w\in W}$ is
multipliable.)

Note: The Lean formalization of this proposition (without sorry) requires the
additional assumption that all FPSs are invertible (i.e., $[x^0]\mathbf{a}_s$
is a unit for each $s$), which is used through Proposition
\ref{prop.fps.prods-mulable-subfams} to guarantee multipliability of subfamilies.
The versions without this invertibility hypothesis are
formalized with sorry.
\end{proposition}

\begin{proof}
This can be derived from the analogous property of finite products, since all
coefficients in a multipliable product are finitely determined (conveniently
using $x^{n}$-approximators, which determine several coefficients at the same
time). See Section \ref{sec.details.gf.prod} for the details.
\end{proof}

\subsection{Fubini rule for infinite products}

\begin{proposition}[Fubini rule for infinite products of FPSs]
\label{prop.fps.prods-mulable-rules.fubini1}
\lean{PowerSeries.tprod_fubini, PowerSeries.tprod_fubini_J, PowerSeries.tprod_fubini_full, PowerSeries.multipliable_tprod_fubini, PowerSeries.multipliable_tprod_fubini_J}
\leantarget
\leanok
Let $I$ and $J$ be two sets. Let $\left(\mathbf{a}_{(i,j)}\right)_{(i,j)\in
I\times J}\in K\left[\left[x\right]\right]^{I\times J}$ be a multipliable
family of FPSs. Assume that for each $i\in I$, the family
$\left(\mathbf{a}_{(i,j)}\right)_{j\in J}$ is multipliable. Assume that for
each $j\in J$, the family $\left(\mathbf{a}_{(i,j)}\right)_{i\in I}$ is
multipliable. Then,
\[
\prod_{i\in I}\;\;\prod_{j\in J}\mathbf{a}_{(i,j)}
=\prod_{(i,j)\in I\times J}\mathbf{a}_{(i,j)}
=\prod_{j\in J}\;\;\prod_{i\in I}\mathbf{a}_{(i,j)}.
\]
(In particular, all the products appearing in this equality are well-defined.)
\end{proposition}

\begin{proof}
The first equality follows by applying Proposition
\ref{prop.fps.prods-mulable-rules.SW1} to $S=I\times J$ and $W=I$ and
$f(i,j)=i$ (and appropriately reindexing the products). The second equality
is analogous. See Section \ref{sec.details.gf.prod} for the details.
\end{proof}

\begin{proposition}[Fubini rule for infinite products, invertible case]
\label{prop.fps.prods-mulable-rules.fubini}
\lean{PowerSeries.fubini_prod_invertible}
\leantarget
\leanok
Let $I$ and $J$ be two sets. Let $\left(\mathbf{a}_{(i,j)}\right)_{(i,j)\in
I\times J}\in K\left[\left[x\right]\right]^{I\times J}$ be a multipliable
family of invertible FPSs. Then,
\[
\prod_{i\in I}\;\;\prod_{j\in J}\mathbf{a}_{(i,j)}
=\prod_{(i,j)\in I\times J}\mathbf{a}_{(i,j)}
=\prod_{j\in J}\;\;\prod_{i\in I}\mathbf{a}_{(i,j)}.
\]
(In particular, all the products appearing in this equality are well-defined.)
\end{proposition}

\begin{proof}
See Section \ref{sec.details.gf.prod}.
\end{proof}
